% =============================================================================
% CAPÍTULO 6 — CONCLUSIONES Y TRABAJOS FUTUROS
% =============================================================================
\chapter{Conclusiones y Trabajos Futuros}
\label{ch:conclusiones}

% TODO: Párrafo introductorio:
%   - Retomar brevemente el problema planteado y los objetivos
%   - Anunciar la estructura del capítulo
\ldots

%% ─────────────────────────────────────────────────────────────
\section{Conclusiones principales}
\label{sec:conclusiones}
%% ─────────────────────────────────────────────────────────────

% TODO: Conclusiones estructuradas por objetivo/dimensión:
%
%   SOBRE EL SISTEMA (OBJ. PRINCIPAL):
%   - Se diseñó e implementó un sistema funcional que asiste la gestión
%     documental para acreditación universitaria (describir qué hace, brevemente)
%   - La arquitectura MAS + LLM + MCP + RAG demostró ser técnicamente viable
%     para el dominio institucional regulado
%   - El sistema cubre el ciclo de vida completo de una propuesta docente:
%     desde la importación hasta la aprobación con trazabilidad
%
%   SOBRE CADA COMPONENTE TÉCNICO:
%   - Extracción: el extractor alcanzó F1 = X sobre las N propuestas
%     del corpus (completar con datos reales del cap. 5)
%   - Validación normativa: el agente detectó correctamente el Y\% de
%     no-conformidades con la Ordenanza 75/23
%   - LLM: las sugerencias generadas fueron valoradas como útiles o muy
%     útiles en el Z\% de los casos por los evaluadores
%   - MCP: el log de auditoría registró el 100\% de los eventos
%     trazables del ciclo de vida (completar)
%
%   SOBRE LA ORIGINALIDAD DE LA CONTRIBUCIÓN:
%   - Primera aplicación de MCP al dominio de educación superior
%   - Primera integración MAS + LLM + MCP para acreditación universitaria
%     bajo el marco normativo CONEAU
%   - El sistema no solo automatiza: habilita trazabilidad auditada
%     que ninguna solución previa provee
\ldots

%% ─────────────────────────────────────────────────────────────
\section{Respuesta a las preguntas de investigación}
\label{sec:respuesta_preguntas}
%% ─────────────────────────────────────────────────────────────

% TODO: Retomar las PG/PE del §3.2 y responder una por una:
%   - PG: ¿Puede un sistema MAS+LLM+MCP asistir la gestión documental
%     con trazabilidad y cumplimiento normativo? → RESPUESTA: ...
%   - PE1: ¿Qué información requiere la CONEAU? → RESPUESTA: ...
%   - PE2: ¿Cómo se estructura la arquitectura multi-agente? → ...
%   - PE3: ¿Qué nivel de precisión de extracción? → ...
%   - PE4: ¿MCP mejora la trazabilidad? → ...
%   Esta sección es corta (1-2 párrafos o tabla de síntesis)
\ldots

%% ─────────────────────────────────────────────────────────────
\section{Contribuciones del trabajo}
\label{sec:contribuciones}
%% ─────────────────────────────────────────────────────────────

% TODO: Listar las contribuciones concretas del trabajo:
%   1. Contribución técnica: el sistema software funcional (artefacto DSR)
%   2. Contribución metodológica: pipeline MAS+LLM+MCP para dominios
%      con requisitos normativos y auditabilidad
%   3. Contribución empírica: evaluación sobre corpus real de propuestas
%      de la carrera \programaTesis de UCASAL
%   4. Contribución conceptual: formalización del flujo de trabajo
%      de acreditación como grafo de estados con eventos auditados
%   5. Contribución práctica: el sistema puede ser desplegado y usado
%      por la institución de origen inmediatamente
\ldots

%% ─────────────────────────────────────────────────────────────
\section{Limitaciones del trabajo}
\label{sec:limitaciones}
%% ─────────────────────────────────────────────────────────────

% TODO: Describir limitaciones honestas:
%   - Corpus de evaluación acotado (N propuestas, 1 carrera, 1 institución)
%   - Extractor rule-based: sensible a variaciones de formato DOCX
%   - Dependencia de API externa (OpenAI): costo, latencia, privacidad
%   - Sin módulo de autenticación completo en la versión evaluada
%   - Evaluación de usabilidad con usuarios reales limitada o pendiente
%   - No contempla (todavía) el proceso de metaevaluación CONEAU completo
%     (la solución cubre la preparación documental, no el proceso de pares)
%   - Escalabilidad no evaluada en producción con múltiples usuarios
\ldots

%% ─────────────────────────────────────────────────────────────
\section{Trabajos futuros}
\label{sec:trabajos_futuros}
%% ─────────────────────────────────────────────────────────────

% TODO: Líneas de trabajo futuro, ordenadas por prioridad/impacto:
%
%   CORTO PLAZO (mejoras directas al sistema actual):
%   1. Extractor ML-based: fine-tuning de un modelo secuencia-a-secuencia
%      (ej. T5, Llama fine-tuned) para mayor robustez ante variaciones de formato
%   2. Autenticación y autorización completa: JWT + OAuth2;
%      gestión de roles institucionales en producción
%   3. Despliegue en la nube: Docker + Railway/AWS/Azure para acceso
%      multi-institución
%   4. Soporte PDF nativo: extracción desde propuestas escaneadas
%      (OCR + layout analysis)
%
%   MEDIANO PLAZO (extensiones del sistema):
%   5. Módulo de análisis comparativo: detectar automáticamente
%      inconsistencias entre propuestas de la misma carrera
%      (coherencia vertical en la currícula)
%   6. Dashboard de gestión curricular: indicadores CONEAU agregados
%      a nivel carrera (porcentaje de conformidad, estado de cada asignatura)
%   7. Integración con SIU-Guaraní u otros SIS institucionales
%   8. Soporte multiinstitucional: MCP servers por institución,
%      coordinados por un Host central
%
%   LARGO PLAZO (investigación futura):
%   9. Protocolo A2A: escalar la arquitectura a colaboración entre
%      agentes de distintas instituciones para benchmarking curricular
%   10. Evaluación a gran escala: replicar el estudio en otras carreras
%       e instituciones para validez externa
%   11. Modelos LLM open-source locales (Llama, Mistral) para eliminar
%       la dependencia de APIs externas y garantizar privacidad total
%   12. Integración con plataformas de gestión académica (Moodle, etc.)
\ldots
